\subsection{H3: killed functions with unchecked calls}

We ran every body-replacement mutant of every eligible function (\HThreeFns{} functions), then \HThreeCalls{} sampled call-statement deletions. \HThreeFnsFullyKilled{} functions had every replacement killed; most of the rest could not be assessed because their body mutants were uncovered, not activated or invalid.

Of the \HThreeSurvivors{} activated, valid call-deletion survivors, \HThreeShare{} lay inside a fully killed function (95\% CI \HThreeLo{} to \HThreeHi{}). H3 is \textbf{\HThreeSupported}. Viewed from the functions' side, \HThreeFKWithSurv{} of the \HThreeFKFns{} fully killed functions that had at least one valid sampled deletion contained a deletion no test noticed (Table~\ref{tab:h3}).

The share is higher than the 48\% reported for Java~\cite{maton2024pseudo}. The two numbers are not directly comparable. We delete only call statements, not all statements, and our function-level population is restricted by the mutmut eligibility rule. Still, the direction is the same and stronger: in these Python projects, most of the effects that tests fail to check sit inside functions that whole-function mutation would report as well tested.

A random sample of 40 survivors (seed 1) was dominated by four kinds of call: \texttt{super().\_\_init\_\_(...)} in exception classes, whose message and position attributes no test checks (six of the 40, all in tomlkit); output side effects such as \texttt{flush()} and progress-bar rendering in click; cleanup and invalidation such as \texttt{\_validation\_cache.clear()}, \texttt{\_reindex()} and a linked-list \texttt{unlink()} in a TTL cache; and \texttt{sort()} calls whose ordering effect no test observes. Whether each matters is a review question. The audit can only show that no test noticed, which is why it reports these as candidate gaps rather than defects.

\begin{table}[t]
\centering\small
\begin{tabular}{lrrrrr}
\toprule
Project & Call deletions & Valid & Survived & In fully killed fn & Share \\
\midrule
attrs & 160 & 107 & 8 & 2 & 0.25 \\
more-itertools & 110 & 100 & 4 & 2 & 0.50 \\
dateutil & 183 & 142 & 19 & 11 & 0.58 \\
toolz & 33 & 32 & 2 & 2 & 1.00 \\
marshmallow & 94 & 77 & 9 & 8 & 0.89 \\
click & 300 & 252 & 37 & 31 & 0.84 \\
tomlkit & 257 & 233 & 48 & 34 & 0.71 \\
cachetools & 114 & 98 & 13 & 10 & 0.77 \\
\bottomrule
\end{tabular}
\caption{H3: sampled call-statement deletions. \emph{In fully killed fn}: survivors whose enclosing function had every body replacement killed.}
\label{tab:h3}
\end{table}

